\documentclass[10pt]{article}

% NeurIPS 2026 workshop proposal -- cleaned source-faithful draft.
% This version preserves the proposal structure and substance from the
% uploaded source, adds explicit Part I / Part II headers, fixes minor
% presentation issues, and wires verified literature/workshop references
% into the prose.
%
% Page budget (NeurIPS 2026 hard limit):
%   - Pages 1-3: main proposal
%   - Pages 4-5: organizer information
%   - Reference materials / appendices (reviewer roster): unlimited.

% THIS IS THE PREAMBLE
\usepackage{booktabs}
\usepackage{tabularx}
\usepackage{array}
\usepackage{xcolor}
\usepackage{float}
\definecolor{AIDaRBlue}{HTML}{0645AD}

\newcolumntype{L}[1]{>{\raggedright\arraybackslash}p{#1}}
\newcolumntype{Y}{>{\raggedright\arraybackslash}X}

\newcommand{\speakername}[1]{\textcolor{AIDaRBlue}{#1}}
\newcommand{\confirmed}{\textit{confirmed}}
\newcommand{\tentative}{\textit{tentative}}

% \IfFileExists{neurips_2026.sty}{%
%   \usepackage[final]{neurips_2026}%
% }{%
%   \usepackage[letterpaper,top=0.75in,bottom=0.85in,left=0.8in,right=0.8in]{geometry}%
%   \usepackage{times}%
% }

\usepackage[nonanonymous, nonatbib,main]{neurips_2026}

\usepackage[utf8]{inputenc}
\usepackage[T1]{fontenc}
\usepackage{hyperref}
\usepackage{url}
\usepackage{booktabs}
\usepackage{amsmath}
\usepackage{amssymb}
\usepackage{microtype}
\usepackage{xcolor}
\usepackage{array}
\usepackage{tabularx}
\usepackage{enumitem}
\usepackage{indentfirst}
% \usepackage{parskip}
\usepackage{titlesec}
\usepackage{multicol}
\usepackage{fancyhdr}

\usepackage{caption}

\usepackage{multicol}
\usepackage{fancyhdr}

\captionsetup[table]{skip=4pt}

% Tighten vertical rhythm so the main proposal fits in 3 pages.
%\setlength{\parskip}{1.8pt}

%\setlength{\parindent}{0pt}
\setlength{\parindent}{2em}
\setlist[itemize]{leftmargin=1.1em,itemsep=0pt,topsep=0pt,parsep=0pt}

% Compact research-paper-style headings. Sections and subsections both display.
\titleformat*{\section}{\normalsize\bfseries}
% \titleformat*{\subsection}{\normalsize\bfseries}
% \titleformat{\subsubsection}[runin]{\normalfont\bfseries\itshape}{}{0pt}{}[\,]
\titlespacing*{\section}{0pt}{6pt}{3pt}
\titlespacing*{\subsection}{0pt}{4pt}{2pt}
% \titlespacing*{\subsubsection}{0pt}{2pt}{0.4em}
\titlespacing*{\paragraph}{0pt}{2pt}{0.4em}

% Numbering depth: number sections and subsections (and show them in the body).
\setcounter{secnumdepth}{3}

% Tighter title block.
% \makeatletter
% \renewcommand{\maketitle}{%
%   \begin{center}%
%     {\Large\bfseries \@title\par}%
%     \vspace{2pt}%
%     {\normalsize \@author\par}%
%   \end{center}%
%   % \vspace{-4pt}%
% }
\makeatother

\hypersetup{
  colorlinks=true,
  linkcolor=black,
  citecolor=black,
  urlcolor=blue
}


% Visible proposal split and running headers.
\newcommand{\currentproposalpart}{Main Workshop Proposal}
\newcommand{\setproposalpart}[1]{\gdef\currentproposalpart{#1}}
\setlength{\headheight}{4pt}
\pagestyle{fancy}
\fancyhf{}
\fancyhead[L]{\footnotesize\currentproposalpart}
\fancyhead[R]{\footnotesize AIDaR @ NeurIPS 2026\quad\thepage}
\renewcommand{\headrulewidth}{0.25pt}
\newcommand{\proposalbanner}[2]{%
  % \vspace{-4pt}%
  \noindent\colorbox{black!8}{\parbox{\dimexpr\linewidth-2\fboxsep\relax}{\textbf{#1}\hfill\textit{#2}}}%
  \vspace{2pt}%
}


\title{AIDaR: AI Data Readiness for Scientific AI}
%\title{AIDaR: Data Infrastructure and Benchmarks for Scientific AI}

\date{}

\begin{document}

\setproposalpart{PART I -- Main Workshop Proposal}
\maketitle

% \begin{center}
% \textit{\textbf{How do we know scientific data is ready for reliable, trustworthy AI?}}
% \end{center}

\vspace{-64pt}
\begin{center}
%\textit{From experimental data to reliable scientific AI systems.}\\
\textit{\textbf{Tagline: }Data Infrastructure and Benchmarks for Reliable Scientific AI Systems.} \url{https://go.roche.com/aidar}
\end{center}



% \vspace{-2pt}

% \begin{abstract}
% Scientific AI is moving from curated prediction tasks toward systems that work directly with experimental data, scientific knowledge, and analysis workflows. Foundation models, retrieval systems, and agents are increasingly expected to query evidence, run analyses, interpret outputs, and support follow-up decisions. For these systems to be reliable, data must remain connected to the samples, assays, protocols, measurements, workflow outputs, and contextual assumptions needed to interpret results beyond the lab or notebook where they were produced.

% \textbf{The workshop centers on one question:}
% \textit{which measurable properties of scientific data ecosystems predict downstream AI behavior?} We use AI Data Readiness (AIDaR) to describe this measurable relationship between data ecosystems and model or agent performance.

% AIDaR focuses on two linked problems. First, how should scientific data be organized so AI systems can use it: from raw measurements and analysis-ready files to relational or graph-structured knowledge, multimodal biomedical measurements, workflow records, and large datasets used for model training and experimental feedback? Second, how should evaluations determine where frontier models are reliable: across biological reasoning, practical data analysis, retrieval over structured scientific knowledge, assay-specific workflows, and deployed scientific AI settings?

% The workshop will bring together researchers, industry scientists, and engineers building scientific data systems, foundation models, agents, biomedical evaluations, and industrial assay platforms. The program includes invited talks, contributed papers, demos, panels, and breakout groups focused on data infrastructure, practical evaluations, and lessons from scientific AI deployments.
% \end{abstract}


% =====================================================================
% 1. A Provisional AI Data Readiness Framework
% =====================================================================

\section{Workshop Themes}

AIDaR is organized around two linked themes: data that scientific AI systems can operate on, and evaluations that test whether systems use that data correctly. Unlike workshops that focus primarily on new models, datasets, or benchmark leaderboards, AIDaR makes the scientific data ecosystem itself the object of evaluation: what context it retains, how it is exposed to AI systems, and how those choices affect downstream behavior.

\textbf{Theme 1: AI-facing scientific data infrastructure.} Scientific data must move from raw measurements to formats usable by models, retrieval systems, and agents while preserving links to samples, assay conditions, instruments, protocols, workflow provenance, uncertainty estimates, and analysis-ready outputs. Topics include infrastructure for scientific foundation models, relational and graph-structured data, multimodal biomedical measurements, materials and chemistry datasets, climate or simulation outputs, workflows, and large industrial datasets used for model training and experimental feedback.

\textbf{Theme 2: Benchmarks and evaluations for scientific AI.} Scientific AI systems should be tested on tasks that resemble real scientific work: choosing inputs, running analyses, checking controls, retrieving structured knowledge, interpreting uncertainty, and deciding whether results are reliable enough for follow-up. Topics include biological reasoning, data-analysis tasks, multimodal biomedical evaluation, practical biology benchmarks, scientific retrieval benchmarks, workflow-level failure analysis, and studies of how performance changes across assays, platforms, simulations, instruments, and data organization choices.

Recent work on FAIR scientific data, dataset documentation, ML-ready metadata, AI-ready data, and scientific data-readiness pipelines motivates this workshop. AIDaR connects those efforts to scientific foundation models, retrieval systems, agents, and practical scientific AI evaluation~\cite{wilkinson2016fair,gebru2021datasheets,pushkarna2022datacards,akhtar2024croissant,brewer2025scientific,brewer2026pipeline,majithia2026actionable,hiniduma2024survey}.

% =====================================================================
% 2. Why Now
% =====================================================================
\section{Why Now}
\label{sec:why-now}

Scientific AI is moving from models trained on curated datasets to systems that interact with experimental data, scientific literature, databases, and lab workflows. These systems need data that preserves sample, assay, protocol, measurement, instrument, and workflow context. Modern scientific data systems are not organized for this style of use. Raw instrument data, sample tables, images, count matrices, simulation outputs, workflow records, notebooks, and plots are often split across storage systems, LIMS/ELN tools, workflow runners, and lab-specific conventions. Large industrial datasets add further pressure: perturbation and multiplexed screens, tissue-model systems, and wet-lab feedback loops must be connected to model-training datasets and analysis tools. Evaluations now shape where frontier models can be used in scientific work. A model that performs well on general reasoning may still fail on a specific assay, platform, instrument, analysis workflow, or data modality. Practical evaluations help identify where agents can assist scientists, where human review is required, and where results are not reliable enough to support follow-up experiments or operational decisions. AIDaR is timely because these problems sit between model development and scientific use. The workshop brings together researchers, industry scientists, and engineers building scientific data systems, foundation models, agents, biomedical evaluations, industrial assay platforms, and deployed analysis environments.

\paragraph{Target audience}
AIDaR is for researchers, scientists, and engineers building AI systems that must work with real scientific data. The audience includes AI-for-science researchers, foundation-model and agent builders, scientific RAG researchers, computational biologists, biomedical ML researchers, benchmark and evaluation researchers, data-management and provenance researchers, and teams responsible for assay data, workflow systems, LIMS/ELN integrations, model-training datasets, and deployed analysis tools. We expect a medium-sized workshop: \textbf{100--200 attendees} and approximately \textbf{40--50 submissions}. We will welcome both methods papers and practitioner reports, including case studies, workflow and tool demos, evaluation suites, failure analyses, assay metadata schemas, retrieval benchmarks, agent task suites, and lessons from deployed systems.

\paragraph{Difference from related workshops} 
\label{sec:relationship}

AIDaR sits between broad AI-for-science workshops, dataset venues, and reliable-ML workshops. AI for Science, ML4PS, and FM4LS emphasize scientific ML methods and domain applications; Data-Centric ML and NeurIPS Datasets \& Benchmarks emphasize datasets, benchmarks, and data-centric methods; Reliable ML from Unreliable Data emphasizes corrupted, shifted, adversarial, or strategic data; and ADSD 2025 at ICDM focuses on creating and curating AI-ready scientific datasets~\cite{ai4science2025,ml4ps2025,fm4ls2025,dmlr2025,reliableml2025,adsd2025}. AIDaR is complementary but distinct: its focus is the scientific data ecosystem that determines how AI systems retrieve, analyze, and act on evidence. Compared with broad AI-for-science workshops, AIDaR is more focused on data infrastructure and evaluation conditions than on new modeling methods alone; compared with dataset venues, it is more focused on use rather than publication; and compared with reliable-ML workshops, it is more focused on scientific provenance, workflow structure, assay and instrument context, and deployed analysis environments. To our knowledge, no recent NeurIPS workshop has made AI-facing scientific data readiness itself the central object of evaluation~\cite{brewer2025scientific, brewer2026pipeline, majithia2026actionable}.
\paragraph{Workshop History and Concurrent Submissions}
\label{sec:history}
This workshop has not previously been held at any venue. No listed organizer is concurrently proposing another NeurIPS 2026 workshop.

% =====================================================================
% 3. Program, speakers, logistics, and compliance
% =====================================================================
\section{Program Format, Speakers, and Diversity}
\label{sec:program}

% \paragraph{Program Format}
\paragraph{Interactive one-day format}
\label{sec:format}
AIDaR is a one-day in-person workshop within the 7--9 hour NeurIPS format. The program avoids a talks-only structure: invited talks are short, contributed work and posters/demos are central.  Breakout groups produce written inputs for the post-workshop report. Remote presentation will be used only for emergencies (within the NeurIPS one-hour cap).

\vspace{-6mm}
\begin{table}[htb!]
\centering
\scriptsize

\caption{Draft workshop agenda.}
\setlength{\tabcolsep}{4pt}
\renewcommand{\arraystretch}{1.12}
\begin{tabular}{@{}p{0.12\linewidth}p{0.4\linewidth}p{0.4\linewidth}@{}}
\toprule
\textbf{Time} & \textbf{Session} & \textbf{Focus / AIDaR lens} \\
\midrule

08:45--09:10 &
\textbf{Opening.} Arrival, question wall, opening remarks, and workshop themes &
\textbf{Framing:} what scientific data must retain for reliable AI use \\

09:10--10:10 &
\textbf{Invited talks I.} Max Welling \textit{(University of Amsterdam)}; 

Jure Leskovec \textit{(Stanford)} &
\textbf{AI-operable data:} foundation models; discovery platforms; relational scientific data; structured retrieval and reasoning \\

10:10--10:30 &
\textbf{Coffee break}  &  \\


10:30--11:15 &
\textbf{Contributed spotlights I.} 

Scientific data systems and practical evaluations &
\textbf{Short talks:} schemas; metadata; workflow interfaces; retrieval tasks; agent benchmarks; failure cases \\

11:15--12:00 &
\textbf{Panel 1.} Benchmarking AI on practical biology &
\textbf{Evaluation design:} biological reasoning; biomedical workflows; realistic data-analysis tasks \\

12:00--13:15 &
\textbf{Poster/demo session 1 + lunch} &
\textbf{Community discussion:} accepted papers; demos; tools; datasets; benchmark proposals \\

13:15--14:00 &
\textbf{Invited talks II.}  Marianna Rapsomaniki \textit{(University of Lausanne)}; 

Chloe-Agathe Azencott \textit{(Mines Paris--PSL)} &
\textbf{Experimental connection:} assay design; sample context; perturbations; multimodal measurements; reproducibility \\

14:00--14:50 &
\textbf{Contributed spotlights II.} 

Analysis histories and deployment lessons &
\textbf{Short talks:} workflow records; human review; failed analyses; reproducible AI-assisted science \\

14:50--15:40 &
\textbf{Poster/demo session 2 + coffee} &
\textbf{Extended discussion:} deployed systems; benchmark tasks; datasets; lessons from practice \\

15:40--16:25 &
\textbf{Breakout groups.} AI-facing data infrastructure &
\textbf{Small-group work:} perturbation datasets; multiplexed screens; tissue-scale systems; training data; wet-lab feedback \\

16:25--17:05 &
\textbf{Panel 2.} Engineering AI-operable assay data systems &
\textbf{Infrastructure:} metadata; raw instrument data; analysis-ready files; provenance; versioning; scientist-facing interfaces \\

17:05--17:45 &
\textbf{Report-out and closing.} Breakout summaries, community roadmap, and post-workshop report plan &
\textbf{Outputs:} roadmap; shared problems; evaluation priorities; post-workshop synthesis \\

\bottomrule
\end{tabular}
\end{table}
\vspace{-2pt}

\paragraph{Invited talks \& confirmed speakers}
\label{sec:speakers}

Max Welling (University of Amsterdam) and Jure Leskovec (Stanford) open with AI-operable scientific data systems: data for foundation models and discovery platforms, and relational/graph infrastructure for retrieval and reasoning. Marianna Rapsomaniki (University of Lausanne) and Chloe-Agathe Azencott (Mines Paris–PSL) then ground the evaluation theme in biomedical data: assay design, sample and tissue context, treatment conditions, multimodal measurements, biological signal, and reproducibility. Confirmed invited speakers have indicated in-person availability, with final venue-specific participation reconfirmed after location assignment.

\paragraph{Panels \& confirmed panelists}
\label{sec:panels}

The panels are designed to open a discussion around the themes of the workshop. Panel 1, ``Benchmarking AI on practical biology,'' will focus on how to construct evaluations that approximate realistic scientific work, especially biological reasoning and data-analysis tasks, rather than only ranking model accuracy on static benchmarks. Confirmed panelists include Harihara Muralidharan (LatchBio), Kexin Huang (Phylo/Biomni), Pablo Meher Rojas (IBM/Dream Challenge), and Jon Laurent (Edison Scientific). Panel 2, ``Engineering AI-operable assay data systems,'' will focus on the infrastructure required to make experimental data usable by scientific AI systems: metadata, raw instrument data, analysis-ready files, workflow data, version history, and APIs. Confirmed panelists include Hoifung Poon (Microsoft Research), Jason Fries (Stanford) and industry panelists from TakaraBio, Vizgen, AtlasXOmics, and Portal (under discussion). 

\paragraph{Breakout groups}
\label{sec:breakouts}
Breakout groups will examine AI-facing data infrastructure for large industrial datasets. We are in discussion with participants from ManifoldBio, Polyphron, Tahoe Therapeutics, and LigandAI, and will finalize leads after acceptance. Groups will discuss how perturbation datasets, multiplexed screens, assay readouts, model-training data, and wet-lab feedback should be organized for model and agent use, then report practical recommendations. Breakouts will produce public problem statements and evaluation-design recommendations, not company-specific plans.


\paragraph{Diversity, Inclusion, and Community Building}
\label{sec:diversity}

AIDaR brings together communities that often work on adjacent problems but rarely meet in one venue: AI-for-science researchers, scientific data-infrastructure builders, benchmark, evaluation, and data-management experts, and practitioners deploying AI in scientific settings. The organizing team and invited speakers span academia, industry research, biotechnology, and data systems, with diversity in geography, seniority, discipline, and research perspective. The workshop is designed to be participatory rather than talks-only. Posters, demos, contributed spotlights, panels, and breakouts will give early-career researchers, practitioners, tool builders, and domain scientists visible ways to contribute. The CFP will welcome early-stage work, tools, datasets, benchmark failures, negative results, deployment lessons, and postmortems of AI failures linked to data conditions. We will use broad CFP and reviewer outreach to reach beyond the organizers' immediate networks and help seed a durable cross-domain community.


\section{Logistics, review process, and compliance}

\paragraph{Submissions, review, COI, and LLM policy}

All contributed work will be managed through OpenReview. We will accept non-archival short papers, extended abstracts, benchmarks, tools, demos, case studies, and position papers. The CFP will state the non-archival status, access policy, and that work already accepted to an archival ML venue is discouraged. Invited speakers will be asked to present content that is novel or substantially tailored to AIDaR. Each submission will receive three reviews, with reviewer load capped at three papers where possible; if submissions exceed expectations, we will recruit additional reviewers or area chairs. Organizers will not give invited talks, submit papers, or participate in decisions for conflicted submissions. Reviewers will be instructed not to use LLMs or other agents, consistent with NeurIPS 2026 workshop policy.



\vspace{-1mm}

\paragraph{Timeline}
\label{sec:timeline}

\textbf{Jul 2026:} Website, CFP, and OpenReview launch. 
\textbf{Aug 29, 2026 AoE:} Submission deadline. 
\textbf{Sep 2026:} Review and meta-review, with notifications by \textbf{Sep 29, 2026 AoE}. 
\textbf{Dec 11--12, 2026:} Sydney workshop dates; \textbf{Dec 12--13, 2026:} Paris/Atlanta workshop dates.

\vspace{-1mm}

\paragraph{Location preference and technical needs} Preferred order: {Paris, Atlanta, Sydney}; any venue acceptable if needed. Hardware needs are standard -- projector/screen, podium, Q\&A and panel microphones/table, poster boards, and  recording/livestream support. 

\vspace{-1mm}
\section{Expected Outcomes}
AIDaR will seed a cross-domain community around AI-facing scientific data readiness: the measurable data conditions that shape whether scientific AI systems can reliably retrieve, analyze, and act on evidence. Outputs will include a public post-workshop report, a taxonomy of readiness failures, reusable data-system patterns for samples, assays, workflows, model-training data, and wet-lab feedback, and practical recommendations for task-specific scientific AI evaluation. The workshop will also define a roadmap for future shared tasks, benchmark proposals, and community evaluation artifacts.



% =====================================================================
% 6. Organizer Information -- two pages maximum (NeurIPS 2026 hard limit).
% =====================================================================
\newpage
\setcounter{section}{0}
\setproposalpart{PART II -- Organizer Information}
\begin{center}
    \textbf{\large Organizer Information}
\end{center}

% \subsection{Note on This Section}
% \label{sec:org-note}
% Per NeurIPS 2026 guidance, this section is limited to two pages.

\section{COI and Concurrent-Proposal Statement}
\label{sec:org-coi}
Each listed organizer has confirmed in writing, prior to OpenReview submission, that (a)~they are not listed on more than two NeurIPS 2026 workshop proposals, (b)~they understand that organizers cannot give invited talks at this workshop and will limit their on-stage role to brief introductions or panel moderation, and (c)~they will not submit papers to this workshop, nor will their students, postdocs, or close collaborators. Organizers will not review or shape acceptance decisions about submissions from within their own organization (corporation-wide for large corporations).

\section{Organizing Team}
\label{sec:org-team}

\begin{itemize}
    \item \textbf{Michal Rosen-Zvi.}
\textit{The Hebrew University of Jerusalem; IBM Research; 
\href{mailto:michal.rosen-zvi@mail.huji.ac.il}{michal.rosen-zvi@mail.huji.ac.il}.}
Director of AI for Healthcare and Life Sciences at IBM Research and Adjunct Professor of Computational Medicine at the Hebrew University, working on machine learning for healthcare, biomedical AI, scientific reasoning, and AI-driven discovery systems.

\item \textbf{Zoe Piran.}
\textit{Stanford University; Roche / Genentech;
\href{mailto:piran.zoe@gene.com}{piran.zoe@gene.com}.}
Postdoctoral Fellow at Genentech and Stanford University working at the intersection of machine learning and computational biology, with research on cellular dynamics from multimodal biological data and scientific discovery workflows.

\item \textbf{Kenny Workman.}
\textit{LatchBio;
\href{mailto:kenny@latch.bio}{kenny@latch.bio}.}
Co-founder and CTO of LatchBio developing infrastructure and benchmarks for data analysis agents. His interests span engineering in biology, computing and AI.

\item \textbf{Edaeni Hamid.}
\textit{Roche / Genentech;
\href{mailto:hamid.edaeni@gene.com}{hamid.edaeni@gene.com}.}
Machine learning and data-platform practitioner focused on deploying AI systems across complex scientific and enterprise data ecosystems, including AI data readiness, scientific data infrastructure, evaluation, and operational AI in research environments.

\item \textbf{Vladimir Ermakov.}
\textit{Edison Scientific;
\href{mailto:vermakov@edisonscientific.com}{vermakov@edisonscientific.com}.}
Researcher and entrepreneur building AI systems for scientific discovery, biomedical research, and foundation-model applications, with work on scientific reasoning, evaluation, evidence synthesis, and hypothesis generation.

\item \textbf{Arindam Sett.}
\textit{Roche / Genentech;
\href{mailto:sett.arindam@gene.com}{sett.arindam@gene.com}.}
Principal Machine Learning Engineer at Genentech working on AI data readiness, scientific AI platforms, retrieval-augmented systems, and agentic workflows, and leading proposal coordination and outreach around the AIDaR workshop theme.

\item \textbf{Sina Booeshaghi.}
\textit{University of California, Berkeley;
\href{mailto:sinab@berkeley.edu}{sinab@berkeley.edu}.}
Postdoctoral Scholar in Bioengineering developing single-cell genomics technologies, with work on robust assay measurements, sequencing-data processing, and reproducible analysis workflows for multimodal genomic data.


\end{itemize}


% \paragraph{Shirley Ho.}
% \textit{New York University; Simons Foundation; 
% \href{mailto:shirley@polymathic.org}{shirley@polymathic.org}.}
% Astrophysicist and machine learning researcher whose work spans cosmology, foundation models for science, and interpretable machine learning. She has pioneered the application of modern deep learning methods in astrophysics and is a leading advocate for AI-driven scientific discovery.



\section{Why This Team Can Deliver the Workshop}
\label{sec:org-fit}
The team spans industry research labs and academia, with combined expertise across applied ML for biomedical and scientific data, foundation models for science, scientific data infrastructure, and trustworthy/data-centric AI. The team intentionally blends senior organizers with newer workshop organizers to support junior involvement while maintaining organizational experience. Outreach is already underway, with anchor speakers confirmed across AI for Science, foundation models, scientific infrastructure, and biomedical AI. 

The organizing team has assembled an initial reviewer pool of 45 researchers spanning academia, biotech, scientific infrastructure, and industry research, providing immediate capacity for workshop review and community outreach. The workshop website, CFP, OpenReview venue, reviewer recruitment plan, and post-workshop artifact plan all have designated owners; OpenReview venue setup, CFP drafting, and reviewer recruitment have explicit owners on the team. The team is intentionally diverse across gender, institution type (industry research, academic, biotech), geography (US, Israel), and career stage.

\section{Diversity of the Organizing Team}
\label{sec:org-diversity}
The organizing team is diverse across gender, institution type, geography, seniority, and discipline. The team spans pharmaceutical research, academic research, scientific infrastructure, biotech platforms, and independent scientific-AI organizations, with members across the US and Israel. The organizers intentionally bridge scientific ML, biomedical AI, data systems, provenance, foundation models for science, and production AI infrastructure. 

The CFP and PC recruitment will explicitly extend beyond the organizers' immediate professional networks through university mailing lists, professional societies, social media, and cross-institutional outreach to broaden representation across geography, career stage, institution type, and research perspective.

\section{Program Committee Plan}
\label{sec:pc}

We are assembling a Program Committee of 40--50 reviewers before the CFP launch. Assuming 40--50 submissions and three reviews per submission, this target keeps the expected reviewer load at or below three papers per reviewer. If submissions exceed 50, we will expand the PC or appoint area chairs to maintain review quality and timely decisions.

The PC is organized around expertise clusters aligned with the workshop's themes: AI-facing scientific data infrastructure; evaluations for scientific AI; workflow systems and deployed analysis tools; biomedical and scientific AI applications; industrial assay platforms; reproducibility, failure analysis, and human review. Each submission will be assigned to reviewers with matching technical expertise, and COI checks will be performed through OpenReview before assignments are finalized. The named roster below reflects the current confirmed reviewer pool.


\section{Current confirmed reviewer pool.}\par\vspace{2pt}

% \noindent{\scriptsize
\setlength{\tabcolsep}{2pt}
\renewcommand{\arraystretch}{1.05}
\begin{tabular}{@{}p{0.49\textwidth}p{0.49\textwidth}@{}}
R01: J. Karin --- HUJI & R23: J. Liu --- Genentech \\

R02: K. Pang --- Stanford & R24: H. M. Khiem --- University of Cincinnati \\

R03: M. Schaefer --- Stanford & R25: M. Patel --- Walmart \\

R04: K. Ponsin --- Rockefeller & R26: O. Bazgir --- Oracle \\

R05: A. Ghandehari --- UC Irvine & R27: J. Rossen --- Harvard \\

R06: A. Gorla --- UCLA & R28: V. Shitov --- Helmholtz \\

R07: A. Shafran --- ETH Zurich & R29: N. Walter --- CISPA \\

R08: S. Nair --- Genentech & R30: A. Fomitchev --- UC Santa Cruz \\

R09: A. Gupta --- Stanford & R31: S. Joshi --- Intuit \\

R10: M. Bereket --- Stanford & R32: S. Kurant --- Technion \\

R11: A. Wu --- Genentech & R33: Y. Lee --- Xaira \\

R12: N. Mudrik --- Johns Hopkins & R34: Y. Shamay --- Technion \\

R13: R. Mintz --- HUJI & R35: R. Littman --- Genentech \\

R14: R. Friedman --- HUJI & R36: E. Peterfreund --- Yale \\

R15: W. Cheng --- Carnegie Mellon & R37: P. Gajare --- UC Irvine \\

R16: Z. Zhu --- UC Berkeley & R38: N. Azran --- HUJI \\

R17: R. Lacombe --- HCVC & R39: M. Danziger --- IBM \\

R18: P. Ogden --- Manifold Bio & R40: M. Osman --- Polyphron \\

R19: Z. Yang --- LatchBio & R41: L. Nuzhna --- Impetus Grants \\

R20: H. Bhasin --- TwentyTwo Bio & R42: J. Wang --- Harvard \\

R21: P. Pao-Huang --- Stanford & R43: S. Alber --- Stanford \\

R22: T. Kitak --- ETH Zürich & R44: E. De Brouwer --- Genentech \\

& R45: B. Rasoulzadeh --- Apple
\end{tabular}
%}

\vspace{2mm}
\textit{\noindent Additional reviewers will be recruited if needed.}

\section{Other Workshops the Organizers are Concurrently Proposing}
\label{sec:org-concurrent}
No listed organizer is concurrently proposing another NeurIPS 2026 workshop. 

% =====================================================================
% References (unlimited; outside the 3+2 page limit).
% =====================================================================
\clearpage
\setproposalpart{References}
\proposalbanner{References}{}
\renewcommand{\refname}{References}
\begin{thebibliography}{99}

\bibitem{wilkinson2016fair}
M. D. Wilkinson, M. Dumontier, I. J. Aalbersberg, G. Appleton, M. Axton, A. Baak, et al.
\newblock The FAIR Guiding Principles for scientific data management and stewardship.
\newblock \emph{Scientific Data}, 3:160018, 2016. \href{https://doi.org/10.1038/sdata.2016.18}{doi:10.1038/sdata.2016.18}.

\bibitem{Gonzalez-Espinozaaiready}
Alfredo Gonzalez-Espinoza
\newblock AI-ready Research Data
\newblock 2026. \href{https://guides.library.cmu.edu/ResearchDataForAI}{https://guides.library.cmu.edu/ResearchDataForAI}

\bibitem{gebru2021datasheets}
T. Gebru, J. Morgenstern, B. Vecchione, J. W. Vaughan, H. Wallach, H. Daum\'{e} III, and K. Crawford.
\newblock Datasheets for datasets.
\newblock \emph{Communications of the ACM}, 64(12):86--92, 2021. \href{https://doi.org/10.1145/3458723}{doi:10.1145/3458723}.

\bibitem{pushkarna2022datacards}
M. Pushkarna, A. Zaldivar, and O. Kjartansson.
\newblock Data Cards: Purposeful and Transparent Dataset Documentation for Responsible AI.
\newblock In \emph{ACM FAccT}, 2022. \href{https://doi.org/10.1145/3531146.3533231}{doi:10.1145/3531146.3533231}.

\bibitem{akhtar2024croissant}
M. Akhtar, O. Benjelloun, C. Conforti, L. Foschini, J. Giner-Miguelez, P. Gijsbers, S. Goswami, N. Jain, M. Karamousadakis, M. Kuchnik, et al.
\newblock Croissant: A Metadata Format for ML-Ready Datasets.
\newblock In \emph{NeurIPS Datasets and Benchmarks Track}, 2024.
\newblock \url{https://arxiv.org/abs/2403.19546}.

\bibitem{hiniduma2024survey}
K. Hiniduma, S. Byna, and J. L. Bez.
\newblock Data Readiness for AI: A 360-Degree Survey.
\newblock arXiv preprint arXiv:2404.05779, 2024.
\newblock \url{https://arxiv.org/abs/2404.05779}.

\bibitem{hiniduma2024aidrin}
K. Hiniduma, S. Byna, J. L. Bez, and R. Madduri.
\newblock AI Data Readiness Inspector (AIDRIN) for Quantitative Assessment of Data Readiness for AI.
\newblock In \emph{Proc. 36th Intl. Conf. on Scientific and Statistical Database Management (SSDBM)}, 2024. \href{https://doi.org/10.1145/3676288.3676296}{doi:10.1145/3676288.3676296}.

\bibitem{brewer2025scientific}
W. Brewer, P. Widener, V. Anantharaj, F. Wang, T. Beck, A. Shankar, and S. Oral.
\newblock Data Readiness for Scientific AI at Scale.
\newblock In \emph{Proc. 54th Intl. Conf. on Parallel Processing Companion (ICPP Companion)}, 2025.
\newblock \url{https://arxiv.org/abs/2507.23018}.

\bibitem{brewer2026pipeline}
W. Brewer, P. Widener, V. Anantharaj, F. Wang, T. Beck, A. Shankar, and S. Oral.
\newblock Data readiness pipeline patterns for scientific AI at scale: Insights from climate, fusion, life sciences, and materials.
\newblock \emph{AI Magazine}, 2026. \href{https://doi.org/10.1002/aaai.70056}{doi:10.1002/aaai.70056}.

\bibitem{majithia2026actionable}
N. Majithia, T. Carey-Wilson, E. Simperl, and N. Shadbolt.
\newblock An actionable framework for AI-ready data.
\newblock \emph{AI Magazine}, 2026. \href{https://doi.org/10.1002/aaai.70054}{doi:10.1002/aaai.70054}.

\bibitem{bhardwaj2024curation}
E. Bhardwaj, H. Gujral, S. Wu, C. Zogheib, T. Maharaj, and C. Becker.
\newblock The State of Data Curation at NeurIPS: An Assessment of Dataset Development Practices in the Datasets and Benchmarks Track.
\newblock In \emph{NeurIPS Datasets and Benchmarks Track}, 2024. \url{https://arxiv.org/abs/2410.22473}.

\bibitem{wu2024management}
Y. Wu, L. Ajmani, S. Longpre, and H. Li.
\newblock A Systematic Review of NeurIPS Dataset Management Practices.
\newblock In \emph{NeurIPS Datasets and Benchmarks Track}, 2024. \url{https://arxiv.org/abs/2411.00266}.

\bibitem{hbr2026taming}
Harvard Business Review Analytic Services.
\newblock \emph{Taming the Complexity of AI Data Readiness}.
\newblock HBR Analytic Services Pulse Survey, sponsored by Cloudera, 2026.
\newblock \url{https://hbr.org/resources/pdfs/comm/cloudera/TamingtheComplexityofAIDataReadiness.pdf}.

\bibitem{ai4science2025}
AI for Science Community.
\newblock The Reach and Limits of AI for Scientific Discovery: NeurIPS 2025 AI for Science Workshop.
\newblock \url{https://ai4sciencecommunity.github.io/neurips25}.

\bibitem{ml4ps2025}
Machine Learning and the Physical Sciences Workshop Organizers.
\newblock Machine Learning and the Physical Sciences, NeurIPS 2025.
\newblock \url{https://ml4physicalsciences.github.io/2025/}.

\bibitem{fm4ls2025}
FM4LS Workshop Organizers.
\newblock NeurIPS 2025 Workshop on Multi-modal Foundation Models and Large Language Models for Life Sciences.
\newblock \url{https://nips2025fm4ls.github.io/}.

\bibitem{dmlr2025}
Data-centric Machine Learning Research Community.
\newblock Data-centric Machine Learning Research Workshop.
\newblock \url{https://dmlr.ai/}.

\bibitem{reliableml2025}
Reliable ML Workshop Organizers.
\newblock Reliable ML from Unreliable Data: NeurIPS 2025 Workshop.
\newblock \url{https://reliablemlworkshop.github.io/}.

\bibitem{adsd2025}
ADSD Workshop Organizers.
\newblock The 1st Workshop on AI-Ready Data for Science Discovery at ICDM 2025.
\newblock \url{https://cnicds.github.io/ICDM2025/}.

\bibitem{afzal2020datareadiness}
S. Afzal, C. Rajmohan, M. Kesarwani, S. Mehta, and H. Patel.
\newblock Data Readiness Report.
\newblock In \emph{IEEE Intl. Conf. on Smart Data Services}, 2020. \url{https://arxiv.org/abs/2010.07213}.

\end{thebibliography}

\end{document}
